%%%%%%%%%%%%%%%%%%%%%%%%%%%%%%%%%%%%%%%%%
% Medium Length Professional CV
% LaTeX Template
% Version 3.0 (December 17, 2022)
%
% This template originates from:
% https://www.LaTeXTemplates.com
%
% Author:
% Vel (vel@latextemplates.com)
%
% Original author:
% Trey Hunner (http://www.treyhunner.com/)
%
% License:
% CC BY-NC-SA 4.0 (https://creativecommons.org/licenses/by-nc-sa/4.0/)
%
%%%%%%%%%%%%%%%%%%%%%%%%%%%%%%%%%%%%%%%%%

%----------------------------------------------------------------------------------------
%	PACKAGES AND OTHER DOCUMENT CONFIGURATIONS
%----------------------------------------------------------------------------------------

\documentclass[
	%a4paper, % Uncomment for A4 paper size (default is US letter)
	11pt, % Default font size, can use 10pt, 11pt or 12pt
]{resume} % Use the resume class

\usepackage{ebgaramond} % Use the EB Garamond font
\usepackage[italian]{babel}
\usepackage[hidelinks]{hyperref}
\hypersetup{
	pdftitle={CV - Pietro Valerio},
	pdfauthor={Pietro Valerio},
	pdfsubject={Curriculum Vitae},
	pdfkeywords={Ingegneria Meccatronica, ROS 2, Controllo Avanzato, IT}
}

%------------------------------------------------

\name{Pietro Valerio} % Your name to appear at the top

% You can use the \address command up to 3 times for 3 different addresses or pieces of contact information
% Any new lines (\\) you use in the \address commands will be converted to symbols, so each address will appear as a single line.

\address{Arcugnano, Vicenza} % Main address


\address{3488341400 \\ pietro.valerio12@gmail.com} % Contact information

%----------------------------------------------------------------------------------------

\begin{document}

%----------------------------------------------------------------------------------------
%	EDUCATION SECTION
%----------------------------------------------------------------------------------------

\begin{rSection}{Formazione}
	
	\textbf{Università di Padova} \hfill \textit{Ottobre 2019 - Settembre 2022} \\ 
	Laurea triennale in Ingegneria Meccatronica \smallskip \\
	Votazione: 96/110
	
	\textbf{Università di Padova} \hfill \textit{Settembre 2022 - Aprile 2025} \\ 
	Laurea magistrale in Ingegneria Meccatronica \smallskip \\
		Votazione: 105/110\\
	Tesi: Migliorare l'interazione uomo robot attraverso la predizione in tempo reale del movimento umano.
	\begin{itemize}

	\item \textbf{Progettazione software}: Sviluppo di una pipeline in ambiente ROS 2 per l'aggregazione di dati da sensori inerziali, finalizzata alla stima accurata della posa dell'operatore.
	
	\item\textbf{ Deep Learning \& Teleoperazione}: Integrazione di modelli di deep learning (sviluppati dal gruppo di ricerca ospitante) per la predizione del movimento in tempo reale, ottimizzando l'invio di comandi a un robot teleoperato tramite logiche predittive.
	
	\item \textbf{Controllo Avanzato}: Implementazione di algoritmi di visual servoing basati sullo stato dell'arte della letteratura scientifica, garantendo un controllo a bassa latenza e bypassando i sistemi standard di pianificazione del moto.
	
	\end{itemize}

	
\end{rSection}

%----------------------------------------------------------------------------------------
%    PROGETTI ACCADEMICI SECTION
%----------------------------------------------------------------------------------------

\begin{rSection}{Progetti Accademici}
	
	\textbf{Laboratorio di Sistemi di Controllo} \hfill \textit{2024} \\
	Modellistica, simulazione e implementazione di algoritmi di controllo per un pendolo inverso in ambiente Matlab. Progettazione di un controllore LQR con stima dello stato, includendo l'analisi di stabilità e la validazione delle prestazioni. Sviluppo di una strategia di \textit{swing-up} per l'innalzamento automatico del pendolo e controllo del motore tramite \textit{load observer}. Validazione del sistema di controllo sia per il mantenimento dell'equilibrio statico (reiezione dei disturbi) sia durante l'inseguimento di specifiche traiettorie di moto della base (\textit{trajectory tracking}).
	
\end{rSection}

%----------------------------------------------------------------------------------------
%	WORK EXPERIENCE SECTION
%----------------------------------------------------------------------------------------

\begin{rSection}{Esperienze Professionali}

	\begin{rSubsection}{Graphicom s.r.l}{Giugno 2025 - Presente}{IT Department}{Vicenza}
		\item Gestione dell'infrastruttura IT aziendale e supporto tecnico.
		\item Analisi dei requisiti e sviluppo di applicazioni software interne per l'ottimizzazione e l'automazione dei processi aziendali.
		\item Sviluppo di un sistema di ticketing digitale per la gestione della manutenzione, integrando cruscotti analitici per il monitoraggio di KPI critici (Response Time, MTTR, MTTF, MTBF, tasso di riapertura).
		\item Progettazione e sviluppo di un'applicazione mobile (Kotlin) per il tracciamento logistico dei prodotti finiti, implementata per certificare i carichi e azzerare le contestazioni dei clienti.
	\end{rSubsection}



\end{rSection}

%----------------------------------------------------------------------------------------
%	TECHNICAL STRENGTHS SECTION
%----------------------------------------------------------------------------------------

\begin{rSection}{Competenze Tecniche}

	\begin{tabular}{@{} >{\bfseries}l @{\hspace{6ex}} l @{}}
		Linguaggi di Programmazione & C, C++, Matlab, Python,Kotlin\\
		Robotica \& Controlli & ROS 2, Visual Servoing, Controllo Avanzato \\
		Database & PostgreSQL \\
		Strumenti & Git, Vim\\
		Lingue & Italiano (Madrelingua), Inglese (B2)
	\end{tabular}

\end{rSection}

%----------------------------------------------------------------------------------------
%	EXAMPLE SECTION
%----------------------------------------------------------------------------------------

%\begin{rSection}{Section Name}

	%Section content\ldots

%\end{rSection}

%----------------------------------------------------------------------------------------

\end{document}
